\documentclass[a0,portrait]{a0poster}
\input{pacotes}
%-------------------------------------------------------
\begin{document}
\onehalfspacing
%-------------------------------------------------------
\input{cabecalho}
\begin{minipage}[b]{1.0\linewidth} %default 0.75
\begin{center}
\Huge \textbf{TÍTULO}\\
\huge \textbf{NOME COMPLETO DO Autor(A)$^{1}$, Autor(A)$^{2}$, Autor(A)$^{3}$, Autor(A)$^{4}$, Autor(A)$^{5}$}\\
%\huge \textbf{}\\
[0.5cm] % Author(s)
\huge Nome da Universidade, UF, Brasil\\[0.4cm] % University/organization
\texttt{email$^{1}$, email$^{2}$, email$^{3}$, email$^{4}$, email$^{5}$}\\
\end{center}
\end{minipage}
{\color{green}{\hrule height 7pt}}
%
\vspace{1cm} % A bit of extra whitespace between the 
\begin{multicols}{2} %
\justify % Justify all text
%--------------------------------------------------------
\section*{RESUMO}
\vspace{-1cm}
O resumo deve ser do \textbf{tipo informativo} e corresponder a uma \textbf{síntese precisa do trabalho}, contendo \textbf{obrigatoriamente} introdução, objetivo, matriz teórica, metodologia, resultados e conclusão.\\

\hspace{-0,5cm}\textbf{Palavras-Chave:} PALAVRAS-CHAVE
%--------------------------------------------------------
\section*{INTRODUÇÃO}
\vspace{-1cm}
Este é o padrão para pôsteres de trabalhos submetidos na Categoria 1 da 21ª SNCT, destinados à divulgação de \textbf{pesquisas em andamento com resultados parciais}.
A \textbf{INTRODUÇÃO} deve contextualizar o objeto de estudo, contendo questões ou problemas de pesquisa, hipóteses (quando houver) e objetivos. Não deve conter ilustrações e tabelas. 

%--------------------------------------------------------
\section*{OBJETIVOS}
\vspace{-1cm}
OBJETIVOS

%--------------------------------------------------------

\section*{METODOLOGIA}
\vspace{-1cm}
Caracterizar o recorte espacial e temporal da pesquisa quando pertinente. Elencar os fundamentos, caracterizar a pesquisa e apresentar os procedimentos metodológicos de acordo com a natureza do trabalho.

%--------------------------------------------------------
\section*{RESULTADOS}
\vspace{-1cm}
Apresentar, analisar e discutir de forma concisa os resultados. Comparar com estudos científicos semelhantes quando pertinente. Relacionar os resultados com os fundamentos teóricos.

O título da ilustração ou tabela deve ser inserido na parte superior, fonte Arial, tamanho 12, alinhamento justificado, espaçamento entre linhas simples, numerados sequencialmente. 

\textbf{Exemplos:}

\begin{figure}[H]
    \flushleft
    \caption{Título da Ilustração.}
    \includegraphics[width=35cm]{SNCT2024.PNG}\\
    \caption*{\footnotesize Fonte: SNCT, 2024.}
    \label{fig:enter-label}
\end{figure}

\begin{table}[H]
     \flushleft
    \caption{Título da Tabela.}
    \begin{tabular}{ccc}
    \hline
    \multicolumn{1}{c}{\textbf{Negrito}} & \textbf{Negrito} & \textbf{Negrito} \\ \hline
    Texto             & Texto      & Texto            \\
    Texto                                  & Texto            & Texto            \\
    Texto                                  & Texto            & Texto      \vspace{0.2cm}    \\
    \hline
    \end{tabular}
    \caption*{\footnotesize Fonte: SNCT, 2024.}
\end{table} \vspace{-0.8cm} 
%--------------------------------------------------------
\section*{CONCLUSÕES}
\vspace{-1cm}
Sintetizar os resultados mais relevantes. Realizar reflexões teóricas que contribuam com o avanço da compreensão do objeto de estudo e campo de investigação. Evidenciar limitações e desafios da pesquisa e do entendimento do objeto de estudo quando pertinente.
Para as citações utilize o comando \verb|\cite{exemplo1}|. Veja como fica \cite{exemplo1}, \cite[p.19]{exemplo2}.
%--------------------------------------------------------
\printbibliography[title = REFERÊNCIAS]

%--------------------------------------------------------
\section*{AGRADECIMENTOS (Opcional)}
\vspace{-1cm}
Inserir, ap\'os as conclus\~oes, de maneira sucinta os agradecimentos. Se o projeto for financiado por alguma ag\^encia de fomento, citar a fonte.
%--------------------------------------------------------

\end{multicols}
\end{document}
